% foundations/kernel_glossary.tex

\chapter*{Knowledge Kernel Vocabulary}
\phantomsection
\addcontentsline{toc}{chapter}{Knowledge Kernel Vocabulary}
\label{chap:kernel-vocabulary}

\newcommand{\kernelterm}[2]{%
\paragraph{#1}
#2\par\vspace{0.7em}
}

The following vocabulary mirrors the canonical terminology of the
approved Knowledge Kernel as used throughout this thesis. It introduces
no additional theory and serves only as a concise terminological
reference.

% ============================================================
\section*{Engineering Concepts}
% ============================================================

\kernelterm{Engineering Concept}
{An enduring concept used to describe the engineering domain
independently of a particular representation or software
implementation.}

\kernelterm{Engineering Object}
{A durable engineering entity whose identity is independent of its
representations, realizations, observations, and computational models.}

\kernelterm{State}
{A structured description of an engineering object under a specified
condition and within a specified context.}

\kernelterm{Representation}
{A derived description used to visualize, communicate, exchange, or
compute engineering information. A representation is not the
engineering object that it represents.}

% ============================================================
\section*{Identity}
% ============================================================

\kernelterm{Engineering Object Identity}
{The identity by which an engineering object remains distinguishable
through changes of representation, realization, observation, and
project context.}

\kernelterm{Identity Composition}
{The formation of engineering identity from stable constituent identity
aspects and their canonical relations.}

\kernelterm{Constructive Identity}
{The identity aspect determined by the constructive definition and
composition of an engineering object.}

\kernelterm{Alignment Identity}
{The engineering identity of an alignment independently of its sampled,
visual, exchange, or metric representations.}

% ============================================================
\section*{Engineering Knowledge}
% ============================================================

\kernelterm{Engineering Observation}
{Source-grounded information concerning an engineering object, state, or
context that may support engineering knowledge.}

\kernelterm{Engineering Knowledge}
{Engineering information admitted for use in interpretation, evaluation,
constraint, comparison, or decision-making.}

\kernelterm{Knowledge Ownership}
{The responsibility boundary assigning engineering knowledge to the
canonical concept that owns its meaning.}

\kernelterm{Truth Mode}
{The epistemic status under which information is treated, such as
observed, derived, assumed, proposed, or decided.}

\kernelterm{Proposal}
{A candidate engineering change or alternative submitted for evaluation
without yet possessing decision authority.}

\kernelterm{Engineering Decision}
{An accepted engineering choice that resolves an evaluated proposal or
engineering alternative.}

% ============================================================
\section*{Evaluation}
% ============================================================

\kernelterm{Residual}
{A quantified deviation used to assess an observation, state, candidate
solution, or engineering alternative.}

\kernelterm{Engineering Constraint}
{A mandatory condition restricting admissible engineering states,
alternatives, or solutions.}

\kernelterm{Preference}
{A non-mandatory engineering criterion used to compare admissible
alternatives.}

\kernelterm{Candidate Solution}
{An engineering state or configuration proposed for evaluation but not
yet accepted as an engineering decision.}

\kernelterm{Delta Operation}
{An operation expressing an intended change between an existing
engineering condition and a proposed condition.}

\kernelterm{Solver Independence}
{The principle that engineering knowledge, residuals, constraints,
preferences, and decisions are not owned by a particular numerical
solver.}

% ============================================================
\section*{Geometry and State}
% ============================================================

\kernelterm{Alignment}
{A durable engineering object whose defining structure is organized
along an intrinsic longitudinal domain.}

\kernelterm{Intrinsic Geometry}
{Geometry described through relations internal to the engineering object
and independently of a particular metric realization.}

\kernelterm{pose2}
{The planar alignment state comprising planar position and heading along
arc-length.}

\kernelterm{pose3}
{The spatial railway-state concept obtained by extending planar
alignment state with the quantities and context required for spatial
interpretation.}

% ============================================================
\section*{Metric and Spatial Realization}
% ============================================================

\kernelterm{Metric Space}
{A space in which distance and related geometric quantities possess
specified metric meaning.}

\kernelterm{Metric Realization}
{The interpretation of intrinsic engineering information within a
specified metric space.}

\kernelterm{Realization Context}
{The contextual information required to perform and interpret a metric
realization.}

\kernelterm{Spatial Operator}
{An operator relating intrinsic, metric, or spatial descriptions of
engineering information.}

\kernelterm{Physical Realization Space}
{The physical space in which an engineering object is materially or
geometrically realized.}

\kernelterm{Engineering Realization}
{The relation between an intended engineering description and its
physically produced or maintained engineering condition.}

\kernelterm{World-to-Track Mapping}
{A spatial operation relating a world-space position to an intrinsic
track-related interpretation.}

\kernelterm{Track-to-World Mapping}
{A spatial operation realizing intrinsic track-related information in
world space.}

% ============================================================
\section*{Persistent Engineering Objects}
% ============================================================

\kernelterm{SpotObject}
{A durable engineering object admitted into SPOT and available for
persistent engineering relations and project-independent use.}

\kernelterm{SpotObject Identity}
{The stable identity by which a SpotObject remains distinguishable
across representations, imports, projects, and realizations.}
